\uuid{E37t}
\titre{Maximum/minimum}
\auteur{bodin ; hanani ; mzari}
\organisation{exo7}
\licence{CC BY-NC-SA 4.0}
\datecreate{2020-01-20}
\niveau{L1}
\module{Analyse}
\chapitre{Continuité, limite et étude de fonctions réelles}
\sousChapitre{Etude de fonctions}
\difficulte{3}
\theme{Maximum, bijection, maximum, minimum}
\exo7qcmid{lille-1.continuite.029}

\contenu{
\begin{question}
\qclassification{123.04}
\qtags{motcle=maximum/minimum}
Soit $I$ un intervalle de $\Rr$ et soit $f : I \to \Rr$ une fonction continue. Soit $J=f(I)$. Quelles sont les assertions vraies ?
\begin{answers}
    \good{$J$ est un intervalle.}

    \bad{Si $I$ est majoré, alors $J$ est majoré.}

    \good{Si $I$ est fermé borné, alors $J$ est fermé borné.}

    \bad{Si $I$ est borné, alors $J$ est borné.}
\end{answers}
\begin{explanations}
Par une fonction continue, l'image d'un intervalle est un intervalle ;
l'image d'un intervalle fermé et borné est un intervalle fermé et borné.
\end{explanations}
\end{question}
}
